\chapter*{Preface: 1994 edition}

\paragraph{Our point of view}  
We believe that calculus can be for our students what it was for Euler
and the Bernoullis: A language and a tool for exploring the whole
fabric of science.  We also believe that much of the mathematical
depth and vitality of calculus lies in these connections to the other
sciences.  The mathematical questions that arise are compelling in
part because the answers matter to other disciplines as well.

The calculus curriculum that this book represents started with a
``clean slate;'' we made no presumptive commitment to any aspect of
the traditional course.  In developing the curriculum, we found it
helpful to spell out our {\bf starting points}, our {\bf curricular
  goals}, our {\bf functional goals}, and our view of the {\bf impact
  of technology}.  Our starting points are a summary of what calculus
is really about.  Our curricular goals are what we aim to convey about
the subject in the course.  Our functional goals describe the
attitudes and behaviors we hope our students will adopt in using
calculus to approach scientific and mathematical questions.  We
emphasize that what is missing from these lists is as significant as
what appears.  In particular, we did {\em not\/} not begin by asking
what parts of the traditional course to include or discard.

\medskip
\noindent
{\bf Starting Points}
\medskip

\begin{minipage}{343pt}

$\hspace{-12pt}\bullet$ Calculus is fundamentally a way of dealing
  with functional relationships that occur in scientific and
  mathematical contexts.  The techniques of calculus must be
  subordinate to an overall view of the underlying questions.

$\hspace{-12pt}\bullet$ Technology radically enlarges the range of
  questions we can explore and the ways we can answer them.  Computers
  and graphing calculators are much more than tools for teaching the
  traditional calculus.
\end{minipage}

\newpage

\noindent
{\bf Starting Points---continued}
\medskip

\begin{minipage}{343pt}
$\hspace{-12pt}\bullet$ The concept of a dynamical system is central
  to science Therefore, differential equations belong at the center of
  calculus, and technology makes this possible {\em at the
    introductory level}.

$\hspace{-12pt}\bullet$ The process of successive approximation is a
  key tool of calculus, even when the outcome of the process---the
  limit---cannot be explicitly given in closed form.
\end{minipage}

\bigskip
\noindent
{\bf Curricular Goals}
\medskip

\begin{minipage}{343pt}

$\hspace{-12pt}\bullet$ Develop calculus in the context of scientific
  and mathematical questions.

$\hspace{-12pt}\bullet$ Treat systems of differential equations as
  fundamental objects of study.

$\hspace{-12pt}\bullet$ Construct and analyze mathematical models.

$\hspace{-12pt}\bullet$ Use the method of successive approximations to
  define and solve problems.

$\hspace{-12pt}\bullet$ Develop geometric visualization with
  hand-drawn and computer graphics.

$\hspace{-12pt}\bullet$ Give numerical methods a more central role.
\end{minipage}

\bigskip
\noindent

\noindent
{\bf Functional Goals}
\medskip

\begin{minipage}{343pt}
$\hspace{-12pt}\bullet$ Encourage collaborative work.

$\hspace{-12pt}\bullet$ Empower students to use calculus as a language
  and a tool.

$\hspace{-12pt}\bullet$ Make students comfortable tackling large,
  messy, ill-defined problems.

$\hspace{-12pt}\bullet$ Foster an experimental attitude towards
  mathematics.

$\hspace{-12pt}\bullet$ Help students appreciate the value of
  approximate solutions.

$\hspace{-12pt}\bullet$ Develop the sense that understanding concepts
  arises out of working on problems, not simply from reading the text
  and imitating its techniques.
\end{minipage}

\medskip

\noindent
{\bf Impact of Technology}
\medskip

\begin{minipage}{343pt}
$\hspace{-12pt}\bullet$ Differential equations can now be solved
  numerically, so they can take their rightful place in the
  introductory calculus course.

$\hspace{-12pt}\bullet$ The ability to handle data and perform many
  computations allows us to explore examples containing more of the
  messiness of real problems.

$\hspace{-12pt}\bullet$ As a consequence, we can now deal with
  credible models, and the role of modelling becomes much more central
  to our subject.
\end{minipage}

\pagebreak

\noindent
{\bf Impact of Technology---continued}
\medskip

\begin{minipage}{343pt}
$\hspace{-12pt}\bullet$ In particular, introductory calculus (and
  linear algebra) now have something more substantial to offer to life
  and social scientists, as well as to physical scientists, engineers
  and mathematicians.

$\hspace{-12pt}\bullet$ The distinction between pure and applied
  mathematics becomes even less clear (or useful) than it may have
  been.
\end{minipage}

\bigskip
By studying the text you can see, quite explicitly, how we have
pursued the curricular goals.  In particular, every one of those goals
is addressed within the very first chapter.  It begins with questions
about describing and analyzing the spread of a contagious disease.  A
model is built, and the model is a system of coupled non-linear
differential equations.  We then begin a numerical assault on those
equations, and the door is opened to a solution by successive
approximations.

 Our implementation of the functional goals is less obvious, but it is
 still evident.  For instance, the text has many more words than the
 traditional calculus book---it is a book to be read.  Also, the
 exercises make unusual demands on students.  Most exercises are not
 just variants of examples that have been worked in the text.  In
 fact, the text has rather few simple ``template'' examples.

%\addtolength{\textheight}{20pt}

\paragraph{Shifts in Emphasis}  
It will also become apparent to you that the text reflects substantial
shifts in emphasis in comparison to the traditional course.  Here are
some of the most striking:
%\addtolength{\tabcolsep}{12pt}
\begin{center}
        \begin{tabular}{c@{\hspace{30pt}}c}
	\hline
        \multicolumn{2}{c}{\sc How the emphasis shifts:}\\ [2pt]
        {\sc increase} &  {\sc decrease} \\ [1pt]
        \hline \\ [-11pt]
        concepts        &  techniques \\ [1pt]
        geometry        &  algebra \\ [1pt]
        graphs          &  formulas \\ [1pt]
        brute force     &  elegance \\ [1pt]
        numerical       &  closed-form \\ [-3pt]
        solutions       &  solutions \\[2pt]
	\hline
        \end{tabular}
\end{center}

Euler's method is a good example of what we mean by ``brute force.''
It is a general method of wide applicability.  Of course when we use
it to solve a differential equation like $y'(t) = t$, we are using a
sledgehammer to crack a peanut.  But at least the sledgehammer {\em
  does\/} work.  Moreover, it works with coconuts (like $y' = y(1 -
y/10))$, and it will just as happily knock down a house (like $y' =
\cos^2(t))$.  Of course, students also see the elegant special methods
that can be invoked to solve $y' = t $ and $y' = y(1 - y/10)$
(separation of variables and partial fractions are discussed in
chapter 11), but they understand that they are fortunate indeed when a
real problem will succumb to these special methods.

\paragraph{Audience}  
Our curriculum is not aimed at a special clientele.  On the contrary,
we think that calculus is one of the great bonds that unifies science,
and all students should have an opportunity to see how the language
and tools of calculus help forge that bond.  We emphasize, though,
that this is not a ``service'' course or calculus ``with
applications,'' but rather a course rich in mathematical ideas that
will serve all students well, including mathematics majors.  The
student population in the first semester course is especially diverse.
In fact, since many students take only one semester, we have aimed to
make the first six chapters stand alone as a reasonably complete
course.  In particular, we have tried to present contexts that would
be more or less broadly accessible.  The emphasis on the physical
sciences is clearly greater in the later chapters; this is deliberate.
By the second semester, our students have gained skill and insight
that allows them to tackle this added complexity.

\paragraph{Handbook for Instructors}  
Working toward our curricular and functional goals has stretched us as
well as our students.  Teaching in this style is substantially
different from the calculus courses most of us have learned from and
taught in the past.  Therefore we have prepared a handbook based on
our experiences and those of colleagues at other schools.  We urge
prospective instructors to consult it.

\paragraph{Origins}  
The Five College Calculus Project has a singular history.  It begins
almost thirty years ago, when the Five Colleges were only Four:
Amherst, Mount Holyoke, Smith, and the large Amherst campus of the
University of Massachusetts.  These four resolved to create a new
institution which would be a site for educational innovation at the
undergraduate level; by 1970, Hampshire College was enrolling students
and enlisting faculty.

Early in their academic careers, Hampshire students grapple with
primary sources in all fields---in economics and ecology, as well as
in history and literature.  And journal articles don't shelter their
readers from home truths: if a mathematical argument is needed, it is
used.  In this way, students in the life and social sciences found,
sometimes to their surprise and dismay, that they needed to know
calculus if they were to master their chosen fields.  However, the
calculus they needed was not, by and large, the calculus that was
actually being taught.  The journal articles dealt directly with the
relation between quantities and their rates of change---in other
words, with differential equations.

Confronted with a clear need, those students asked for help.  By the
mid-1970s, Michael Sutherland and Kenneth Hoffman were teaching a
course for those students.  The core of the course was calculus, but
calculus as it is {\em used\/} in contemporary science.  Mathematical
ideas and techniques grew out of scientific questions.  Given a
process, students had to recast it as a model; most often, the model
was a set of differential equations.  To solve the differential
equations, they used numerical methods implemented on a computer.

The course evolved and prospered quietly at Hampshire.  More than a
decade passed before several of us at the other four institutions paid
some attention to it.  We liked its fundamental premise, that
differential equations belong at the center of calculus.  What
astounded us, though, was the revelation that differential equations
could really {\em be\/} at the center---thanks to the use of
computers.

This book is the result of our efforts to translate the Hampshire
course for a wider audience.  The typical student in calculus has not
been driven to study calculus in order to come to grips with his or
her own scientific questions---as those pioneering students had.  If
calculus is to emerge organically in the minds of the larger student
population, a way must be found to involve that population in a
spectrum of scientific and mathematical questions.  Hence, calculus
{\em in context}.  Moreover, those contexts must be understandable to
students with no special scientific training, and the mathematical
issues they raise must lead to the central ideas of the calculus---to
differential equations, in fact.

Coincidentally, the country turned its attention to the undergraduate
science curriculum, and it focused on the calculus course.  The
National Science Foundation created a program to support calculus
curriculum development.  To carry out our plans we requested funds for
a five-year project; we were fortunate to receive the only multi-year
curriculum development grant awarded in the first year of the NSF
program.  This text is the outcome of our effort.

%%%%%%%%%%%%%%%%%%%%%%%%%%%%%%%%%%%%%%%%%%%%%%%%%%%%%%%%%%%%%%%%%%%%%%%%

\section*{Acknowledgements}

Certainly this book would have been possible without the support of
the National Science Foundation and of Five Colleges, Inc.  We
particularly want to thank Louise Raphael who, as the first director
of the calculus program at the National Science Foundation, had faith
in us and recognized the value of what had already been accomplished
at Hampshire College when we began our work.  Five College
Coordinators Conn Nugent and Lorna Peterson supported and encouraged
our efforts, and Five College treasurer and business manager Jean
Stabell has assisted us in countless ways throughout the project.

We are very grateful to the members of our Advisory Board: to Peter
Lax, for his faith in us and his early help in organizing and chairing
the Board; to Solomon Garfunkel, for his advice on politics and
publishing; to Barry Simon, for using our text and giving us his
thoughtful and imaginative suggestions for improving it; to Gilbert
Strang, for his support of a radical venture; to John Truxal, for his
detailed commentaries and insights into the world of engineering.

Among our colleagues, James Henle of Smith College deserves special
thanks.  Besides his many contributions to our discussions of
curriculum and pedagogy, he developed the computer programs that have
been so valuable for our teaching: Graph, Slinky, Superslinky, and
Tint.  Jeff Gelbard and Fred Henle ably extended Jim's programs to the
MacIntosh and to DOS Windows and X Windows.  All of this software is
available on anonymous ftp at emmy.smith.edu.  Mark Peterson, Robert
Weaver, and David Cox also developed software that has been used by
our students.

Several of our colleagues made substantial contributions to our
frequent editorial conferences and helped with the writing of early
drafts.  We offer thanks to David Cohen, Robert Currier and James
Henle at Smith; David Kelly at Hampshire; and Frank Wattenberg at the
University of Massachusetts.  Mary Beck, who is now at the University
of Virginia, gave heaps of encouragement and good advice as a
co-teacher of the earliest version of the course at Smith.  Anne
Kaufmann, an Ada Comstock Scholar at Smith, assisted us with extensive
editorial reviews from the student perspective.

Two of the most significant new contributions to this edition are the
appendix for graphing calculators and a complete set of solutions to
all the exercises.  From the time he first became aware of our
project, Benjamin Levy has been telling us how easy and natural it
would be to adapt our Basic programs for graphing calculators.  He has
always used them when he taught {\em Calculus in Context}, and he
created the appendix which contains translations of our programs for
most of the graphing calculators in common use today.  Lisa Hodsdon,
Diane Jamrog, and Marcia Lazo have worked long hours over an entire
summer to solve all the exercises and to prepare the results as
\LaTeX{} documents for inclusion in the Handbook for Instructors.  We
think both these contributions do much to make the course more useful
to a wider audience.

We appreciate the contributions of our colleagues who participated in
numerous debriefing sessions at semester's end and gave us comments on
the evolving text.  We thank George Cobb, Giuliana Davidoff, Alan
Durfee, Janice Gifford, Mark Peterson, Margaret Robinson, and Robert
Weaver at Mount Holyoke; Michael Albertson, Ruth Haas, Mary Murphy,
Marjorie Senechal, Patricia Sipe, and Gerard Vinel at Smith.  We
learned, too, from the reactions of our colleagues in other
disciplines who participated in faculty workshops on Calculus in
Context.

We profited a great deal from the comments and reactions of early
users of the text.  We extend our thanks to Marian Barry at Aquinas
College, Peter Dolan and Mark Halsey at Bard College, Donald Goldberg
and his colleagues at Occidental College, Benjamin Levy at Beverly
High School, Joan Reinthaler at Sidwell Friends School, Keith Stroyan
at the University of Iowa, and Paul Zorn at St. Olaf College.  Later
users who have helped us are Judith Grabiner and Jim Hoste at Pitzer
College; Allen Killpatrick, Mary Scherer, and Janet Beery at the
University of Redlands; and Barry Simon at Caltech.

Dissemination grants from the NSF have funded regional workshops for
faculty planning to adopt Calculus in Context.  We are grateful to
Donald Goldberg, Marian Barry, Janet Beery, and to Henry Warchall of
the University of North Texas for coordinating workshops.

We owe a special debt to our students over the years, especially those
who assisted us in teaching, but also those who gave us the benefit of
their thoughtful reactions to the course and the text.  Seeing what
they were learning encouraged us at every step.
 
We continue to find it remarkable that our text is to be published the
way we want it, not softened or ground down under the pressure of
anonymous reviewers seeking a return to the mean.  All of this is due
to the bold and generous stance of W. H. Freeman.  Robert Biewen, its
president, understands---more than we could ever hope---what we are
trying to do, and he has given us his unstinting support.  Our
aquisitions editors, Jeremiah Lyons and Holly Hodder, have inspired us
with their passionate conviction that our book has something new and
valuable to offer science education.  Christine Hastings, our
production editor, has shown heroic patience and grace in shaping the
book itself against our often contrary views.  We thank them all.

%%%%%%%%%%%%%%%%%%%%%%%%%%%%%%%%%%%%%%%%%%%%%%%%%%%%%%%%%%%%%%%%%%%%%%%%%

\section*{To the Student}

In a typical high school math text, each section has a ``technique'' which you practice in a series of exercises very like the examples in the text.  This book is different.  In this course you will be learning to use calculus both as a tool and as a language in which you can think coherently about the problems you will be studying.  As with any other language, a certain amount of time will need to be spent learning and practicing the formal rules.  For instance, the conjugation of {\em \^etre\/} must be almost second nature to you if you are to be able to read a novel---or even a newspaper---in French.  In calculus, too, there are a number of manipulations which must become automatic so that you can focus clearly on the content of what is being said.  It is important to realize, however, that becoming good at these manipulations is not the goal of learning calculus any more than becoming good at declensions and conjugations is the goal of learning French.

Up to now, most of the problems you have met in math classes have had definite answers such as ``17,'' or ``the circle with radius 1.75 and center at (2,3).''  Such definite answers are satisfying (and even comforting).  However, many interesting and important questions, like ``How far is it to the planet Pluto,'' or ``How many people are there with sickle-cell anemia,'' or ``What are the solutions to the equation $x^5 + x + 1 = 0$'' can't be answered exactly.  Instead, we have ways to {\bf approximate} the answers, and the more time and/or money we are willing to expend, the better our approximations may be.  While many calculus problems do have exact answers, such problems often tend to be special or atypical in some way.  Therefore, while you will be learning how to deal with these ``nice'' problems, you will also be developing ways of making good approximations to the solutions of the less well-behaved (and more common!) problems.

The computer or the graphing calculator is a tool that that you will need for this course, along with a clear head and a willing hand.  We don't assume that you know anything about this technology ahead of time.  Everything necessary is covered completely as we go along.

You can't learn mathematics simply by reading or watching others.  The only way you can internalize the material is to work on problems yourself.  It is by grappling with the problems that you will come to see what it is you do understand, and to see where your understanding is incomplete or fuzzy.

One of the most important intellectual skills you can develop is that of exploring questions on your own.  Don't simply shut your mind down when you come to the end of an assigned problem.  These problems have been designed not so much to capture the essence of calculus as to prod your thinking, to get you wondering about the concepts being explored.  See if you can think up and answer variations on the problem.  Does the problem suggest other questions?  The ability to ask good questions of your own is at least as important as being able to answer questions posed by others.

We encourage you to work with others on the exercises.  Two or three of you of roughly equal ability working on a problem will often accomplish much more than would any of you working alone.  You will stimulate one another's imaginations, combine differing insights into a greater whole, and keep up each other's spirits in the frustrating times.  This is particularly effective if you first spend time individually working on the material.  Many students find it helpful to schedule a regular time to get together to work on problems.

 Above all, take time to pause and admire the beauty and power of what you are learning.  Aside from its utility, calculus is one of the most elegant and richly structured creations of the human mind and deserves to be profoundly admired on those grounds alone.  Enjoy!

